% GENERATED FILE. Edit canonical lessons, then run npm run generate:books.
% canonical-source-hash: fnv1a64:fc7d00cdec07994b
% canonical-lessons: MR-C138-matar, MR-C138-dalimb, MR-C138-kalingad, MR-C138-naral, MR-C138-pith

\chapter{Peas, Pomegranate, Watermelon, Coconut, Flour}
\label{ch:mr-peas-pomegranate-watermelon-coconut-flour}
\hlchaptermodality{138}

\begin{chapteropening}
\textbf{By the end of this chapter:} \emph{I can say peas, a pomegranate, a watermelon, a coconut and flour.}

Complete the last lesson of chapter 138: I can say peas, a pomegranate, a watermelon, a coconut and flour.
\end{chapteropening}

\section[\texorpdfstring{\mr{मटार} (maṭār)}{मटार (maṭār)}]{\mr{मटार} (maṭār) — peas}

\label{lesson:MR-C138-matar}



\begin{quote}
Before the new one: say the Marathi for garlic, then the Marathi for ginger.
\end{quote}

\subsection*{You'll want to know: \mr{मटार}}
\textbf{\mr{मटार}} — \emph{maṭār} — \textquotedblleft{}peas\textquotedblright{}.

\begin{cousinweb}
\textbf{\mr{मटार}} (\emph{maṭār}) matches Hindi \emph{maṭar}, \textquotedblleft{}pea\textquotedblright{}, which CDIAL traces to a reconstructed \emph{*maṭṭara}, \textquotedblleft{}pea\textquotedblright{}. CDIAL lists Marathi's own form, \emph{maṭh}, \textquotedblleft{}pea\textquotedblright{}, under its variant \emph{*maṭṭha} (CDIAL 9724).
\end{cousinweb}

\begin{grammarlens}[title={What you've built}]
One more word for this chapter.
\end{grammarlens}

\subsection*{Your turn}
\begin{itemize}
\raggedright
  \item \emph{Say it:} \emph{maṭār}
  \item \emph{Say it:} \emph{maṭār}, once more
  \item \emph{Say it:} say \emph{maṭār}
  \item \emph{Recall:} say the Marathi for garlic, then the Marathi for ginger, then say \emph{maṭār} again
\end{itemize}

\begin{tcolorbox}[breakable,colback=teal!4,colframe=teal!35!black,title={Before you move on}]
What does \mr{मटार} mean? (\textquotedblleft{}Peas\textquotedblright{}.) Say it once more.
\end{tcolorbox}
\section[\texorpdfstring{\mr{डाळिंब} (ḍāḷimb)}{डाळिंब (ḍāḷimb)}]{\mr{डाळिंब} (ḍāḷimb) — a pomegranate}

\label{lesson:MR-C138-dalimb}



\begin{quote}
Before the new one: say the Marathi for ginger, then the Marathi for peas.
\end{quote}

\subsection*{You'll want to know: \mr{डाळिंब}}
\textbf{\mr{डाळिंब}} — \emph{ḍāḷimb} — \textquotedblleft{}a pomegranate\textquotedblright{}.

\begin{cousinweb}
\textbf{\mr{डाळिंब}} (\emph{ḍāḷimb}) is a tadbhava, a word inherited from Sanskrit and reshaped by sound change. It comes from \emph{dāḍimba}, \textquotedblleft{}pomegranate\textquotedblright{} (CDIAL 6254). CDIAL gives \emph{dāḍimba} as a variant of Sanskrit \emph{dāḍima}.
\end{cousinweb}

\begin{grammarlens}[title={What you've built}]
One more word for this chapter.
\end{grammarlens}

\subsection*{Your turn}
\begin{itemize}
\raggedright
  \item \emph{Say it:} \emph{ḍāḷimb}
  \item \emph{Say it:} \emph{ḍāḷimb}, once more
  \item \emph{Say it:} say \emph{ḍāḷimb}
  \item \emph{Recall:} say the Marathi for ginger, then the Marathi for peas, then say \emph{ḍāḷimb} again
\end{itemize}

\begin{tcolorbox}[breakable,colback=teal!4,colframe=teal!35!black,title={Before you move on}]
What does \mr{डाळिंब} mean? (\textquotedblleft{}A pomegranate\textquotedblright{}.) Say it once more.
\end{tcolorbox}
\section[\texorpdfstring{\mr{कलिंगड} (kaliṅgaḍ)}{कलिंगड (kaliṅgaḍ)}]{\mr{कलिंगड} (kaliṅgaḍ) — a watermelon}

\label{lesson:MR-C138-kalingad}



\begin{quote}
Before the new one: say the Marathi for peas, then the Marathi for a pomegranate.
\end{quote}

\subsection*{You'll want to know: \mr{कलिंगड}}
\textbf{\mr{कलिंगड}} — \emph{kaliṅgaḍ} — \textquotedblleft{}a watermelon\textquotedblright{}.

\begin{cousinweb}
\textbf{\mr{कलिंगड}} (\emph{kaliṅgaḍ}) is a tadbhava, a word inherited from Sanskrit and reshaped by sound change. It comes from \emph{kāliṅga}, \textquotedblleft{}watermelon\textquotedblright{} (CDIAL 3100). CDIAL gives \emph{kāliṅga} as a variant of Sanskrit \emph{kālinda}, \textquotedblleft{}watermelon\textquotedblright{}.
\end{cousinweb}

\begin{grammarlens}[title={What you've built}]
One more word for this chapter.
\end{grammarlens}

\subsection*{Your turn}
\begin{itemize}
\raggedright
  \item \emph{Say it:} \emph{kaliṅgaḍ}
  \item \emph{Say it:} \emph{kaliṅgaḍ}, once more
  \item \emph{Say it:} say \emph{kaliṅgaḍ}
  \item \emph{Recall:} say the Marathi for peas, then the Marathi for a pomegranate, then say \emph{kaliṅgaḍ} again
\end{itemize}

\begin{tcolorbox}[breakable,colback=teal!4,colframe=teal!35!black,title={Before you move on}]
What does \mr{कलिंगड} mean? (\textquotedblleft{}A watermelon\textquotedblright{}.) Say it once more.
\end{tcolorbox}
\section[\texorpdfstring{\mr{नारळ} (nāraḷ)}{नारळ (nāraḷ)}]{\mr{नारळ} (nāraḷ) — a coconut}

\label{lesson:MR-C138-naral}



\begin{quote}
Before the new one: say the Marathi for a pomegranate, then the Marathi for a watermelon.
\end{quote}

\subsection*{You'll want to know: \mr{नारळ}}
\textbf{\mr{नारळ}} — \emph{nāraḷ} — \textquotedblleft{}a coconut\textquotedblright{}.

\begin{cousinweb}
\textbf{\mr{नारळ}} (\emph{nāraḷ}) is a tadbhava of Sanskrit \emph{nārikela}, \textquotedblleft{}coconut palm and fruit\textquotedblright{} (CDIAL 7075). CDIAL follows Bloch in taking the Sanskrit word from a Dravidian compound, comparing Tamil \emph{nāri}, \textquotedblleft{}fibre\textquotedblright{}.
\end{cousinweb}

\begin{grammarlens}[title={What you've built}]
One more word for this chapter.
\end{grammarlens}

\subsection*{Your turn}
\begin{itemize}
\raggedright
  \item \emph{Say it:} \emph{nāraḷ}
  \item \emph{Say it:} \emph{nāraḷ}, once more
  \item \emph{Say it:} say \emph{nāraḷ}
  \item \emph{Recall:} say the Marathi for a pomegranate, then the Marathi for a watermelon, then say \emph{nāraḷ} again
\end{itemize}

\begin{tcolorbox}[breakable,colback=teal!4,colframe=teal!35!black,title={Before you move on}]
What does \mr{नारळ} mean? (\textquotedblleft{}A coconut\textquotedblright{}.) Say it once more.
\end{tcolorbox}
\section[\texorpdfstring{\mr{पीठ} (pīṭh)}{पीठ (pīṭh)}]{\mr{पीठ} (pīṭh) — flour}

\label{lesson:MR-C138-pith}



\begin{quote}
Before the new one: say the Marathi for a watermelon, then the Marathi for a coconut.
\end{quote}

\subsection*{You'll want to know: \mr{पीठ}}
\textbf{\mr{पीठ}} — \emph{pīṭh} — \textquotedblleft{}flour\textquotedblright{}.

\begin{cousinweb}
\textbf{\mr{पीठ}} (\emph{pīṭh}) is a tadbhava, a word inherited from Sanskrit and reshaped by sound change. It comes from \emph{piṣṭa}, \textquotedblleft{}crushed, ground; flour\textquotedblright{} (CDIAL 8218). It is built on \emph{piṣ}, \textquotedblleft{}to grind\textquotedblright{}.
\end{cousinweb}

\begin{grammarlens}[title={What you've built}]
One more word for this chapter.
\end{grammarlens}

\subsection*{Your turn}
\begin{itemize}
\raggedright
  \item \emph{Say it:} \emph{pīṭh}
  \item \emph{Say it:} \emph{pīṭh}, once more
  \item \emph{Say it:} say \emph{pīṭh}
  \item \emph{Recall:} say the Marathi for a watermelon, then the Marathi for a coconut, then say \emph{pīṭh} again
\end{itemize}

\begin{tcolorbox}[breakable,colback=teal!4,colframe=teal!35!black,title={Before you move on}]
What does \mr{पीठ} mean? (\textquotedblleft{}Flour\textquotedblright{}.) Say it once more.
\end{tcolorbox}
